% Generated by task fix-fragments from dossiers/register.md (rows R12, R57, R105 to R109, R111,
% R112, R116 to R118, R133, and R38 for what is not to be reported), with the finding boxes those
% rows point to. No \section line: 11-options.tex provides it.
\label{sec:gen-report-upstream}

\begin{sloppypar}
Thirteen matters, each a disagreement inside leanVM or a defect of one of its sources, none a
defect of the blueprint. Lines are at leanVM's pin \code{a386121f}: specification files under
\file{doc/leanvm/body/}, Rust files under \file{crates/}, the Python verifier
\file{python-verifier/verifier.py} (\code{py:}), the recursion guest
\file{crates/rec_aggregation/guests/aggregate.py}. The first two change what a verifier accepts
or what an honest prover produces; the next seven are gaps or loose statements of the
specification and the Rust's parameter accounting; the last four are stale comments in the Rust.
\end{sloppypar}

\begin{sloppypar}
\begin{enumerate}[leftmargin=*, itemsep=5pt]
\item \textbf{The public-input check differs between the specification and its three
  verifiers.} The specification checks each limb of the two claimed evaluations against the
  public words, $c_ℓ = (1 + r)\,w_{0,ℓ} + r\,w_{1,ℓ}$ for $ℓ ∈ \{0, 1\}$; the Rust verifier, the
  Python verifier and the recursion guest check one combined equation,
  $c_0 + y\, c_1 = (1 + r)\, w_0 + r\, w_1$ ($y$ the generator whose powers carry a word's limbs), and so accept strictly more transcripts (both are sound,
  the combined one by the review's argument on paper). \emph{Where:} \code{08-end-to-end-protocol.tex:29-33} against
  \code{lean_vm/src/cpu/mod.rs:750-755}, \code{py:1398-1400}, \code{aggregate.py:1680-1683}.
  \emph{Whose:} a disagreement between the specification and all three executable verifiers.
  \emph{Asked:} a clarification: which check is normative, and the specification's §8.2 rewritten
  to it with its error (\code{R12}).

\item \textbf{The Rust prover of Flock is not perfectly complete.} The prover derives a
  coefficient through $(1 + r_{\mathrm{eq}})^{-1}$, and at $r_{\mathrm{eq}} = 1$ it emits a proof
  its own verifier rejects (probability at most $(k_{\mathrm{batch}} + 1)/\abs{\E}$ per proof); the
  specification's protocol is perfectly complete. \emph{Where:}
  \code{flock/src/zerocheck.rs:116-118}. \emph{Whose:} the Rust prover. \emph{Asked:} a fix (the
  prover computes $G(0)$ from its witness rather than solving for it)
  (\code{R111}).

\item \textbf{Seven disagreements in the setup, the commitment and the bus phase}, on which the
  two verifiers agree with each other (\cref{tab:gen-bus-disagreements}). (a) The rate is announced
  with the sizes, the specification naming seven announced values (\code{lean_vm/src/cpu/mod.rs:123,
  149} against \code{08-end-to-end-protocol.tex:55}); (b) the specification's prose names the count
  root before the bus root, the stream has the bus root first (\code{08-end-to-end-protocol.tex:68}
  against \code{lean_vm/src/gkr.rs:272-276}); (c) the seed hashes one constant naming the circuit,
  where the specification says the R1CS matrices (\code{lean_vm/src/cpu/mod.rs:66-93} against
  \code{08-end-to-end-protocol.tex:55}); (d) the specification does not mention seven checks both
  verifiers make (checks 2, 7, 8, 9, 10, 16 and 17 of \cref{tab:gen-bus-checks}); (e) the verifiers
  draw one combiner more than GKR layers, the specification one per layer
  (\code{05-arithmetization.tex:91}; \code{lean_vm/src/gkr.rs:423}); (f) the Python verifier takes
  the program as the stacked table and does not check that its unused slots are zero
  (\code{py:566, 857}); (g) the fingerprint bound is $5·2^μ$ in the Rust and $4·2^μ$ in the
  specification (\code{lean_vm/src/leaf.rs:110} against \code{05-arithmetization.tex:37}).
  \emph{Whose:} the specification for (a) to (e) and (g), which is silent or in another order; the
  Python for (f). \emph{Asked:} clarifications of the specification, and for (f) a check or a
  sentence saying that the table is the statement (\code{R105}).

\item \textbf{The Fiat--Shamir seed binds a constant nobody can recompute at the pin.} The Rust
  binds \code{R1CS_DIGEST}, whose recipe needs matrices the module no longer builds (\enquote{To
  recompute it, check out the last commit that still had \code{build_matrices}}); the Python and
  the recursion guest copy the bytes. \emph{Where:} \code{flock/src/hash.rs:259-275},
  \code{py:629}; the specification's \code{08-end-to-end-protocol.tex:55}. \emph{Whose:} the Rust,
  mirrored by the Python and the guest. \emph{Asked:} a documented constant: the derivation stated
  in the specification, or a recipe that recomputes it at the current revision (\code{R106}).

\item \textbf{The specification gives no round-by-round analysis of the bus phase.} Its one
  round-by-round theorem is the opening's; the GKR's and the fingerprint's per-challenge errors are
  the blueprint's own analysis, and the Rust asserts only a coarse sum. \emph{Where:}
  \code{b-polynomial-commitment-scheme.tex:139-152} (Theorem B.7); \code{lean_vm/src/leaf.rs:108-118}.
  \emph{Whose:} the specification (an omission). \emph{Asked:} a clarification: the errors of the
  bus phase, challenge by challenge (\code{R107}).

\item \textbf{The table sumcheck's batching error is loose by one, and two things are unsaid.}
  The specification charges $(ν_{\mathrm{side}} + B)/\abs{\E}$ where its own Corollary 3.9 gives
  one less; it gives no error for the recycled zerocheck point; and it does not say which
  coefficient of the round polynomial is dropped. \emph{Where:} \code{05-arithmetization.tex:132,
  155}; \code{03-proving-primitives.tex:67} (\code{cor:idtest}), \code{:110}. \emph{Whose:} the
  specification. \emph{Asked:} a fix of the error and two clarifications (\code{R108}).

\item \textbf{The Rust's security accounting unions only the algebraic checks over the Johnson
  list.} The bus's fingerprint term is not multiplied by the list size; at these sizes there are
  about 23 bits of room, so the claim stands, but the accounting does not say so. \emph{Where:}
  \code{pcs/src/whir_config.rs:537-568}; \code{lean_vm/src/leaf.rs:108-118}. \emph{Whose:} the
  Rust's parameter accounting. \emph{Asked:} a clarification in the accounting, or the term
  included (\code{R117}).

\item \textbf{Flock's circuit, layout and constants are fixed only by the Rust and the Python.}
  The BLAKE2s circuit and its wire positions, the slot of each limb, the selector of a limb claim,
  $g_0$, the embedding $φ_8$ and the modulus of $\F_{2^8}$, $k_{\mathrm{batch}} = τ_{\mathrm{BLAKE2S}}$
  and the floor $τ_{\mathrm{BLAKE2S}} ≥ 3$ are in no annex of the specification. \emph{Where:} for
  example \code{flock/src/hash.rs:40-64}, \code{lean_vm/src/hash_flock.rs:87-115},
  \code{pcs/src/stack_open.rs:84-97}, \code{lean_vm/src/cpu/mod.rs:162-166}, \code{py:1092-1095};
  the full list is \cref{f:flock-ring-written-from-spec}. \emph{Whose:} the specification (it is
  incomplete). \emph{Asked:} documented constants, in the specification's Annex C (\code{R57}).

\item \textbf{The witness stack's size has neither floor, ceiling nor tie rule in the
  specification.} The specification's $M = ⌈\log_2 N⌉$ has no bounds, while both verifiers floor
  it at 15 and reject it above 28, and it gives no order for blocks of equal size, which both
  verifiers break by column index. \emph{Where:} \code{04-committing-the-witness.tex:10} against
  \code{lean_vm/src/witness.rs:97}, \code{py:890}, \code{lean_vm/src/cpu/mod.rs:174-176},
  \code{py:1379}. \emph{Whose:} the specification. \emph{Asked:} a clarification (\code{R133}).

\item \textbf{Stale comments on the table sumcheck's round message.} They say the round
  polynomial is sent whole at four nodes, that the verifier checks $h(0) + h(1)$ and interpolates,
  and that $h(0)$ is derived; the code sends three coefficients, derives $c_1$ and checks nothing
  per round. \emph{Where:} \code{lean_vm/src/constraints.rs:24, 26, 257-260, 265-266} against
  \code{fiat_shamir/src/transcript.rs:289-302}. \emph{Whose:} the Rust (comments). \emph{Asked:}
  a fix (\code{R109}).

\item \textbf{Stale comments on what binds Flock's counter and flags.} They say the bytecode
  binds them; the table reads the metadata cell from memory, as the specification says; two more
  comments of the kind misdescribe the Flock round message and the number of ring-switched
  claims. \emph{Where:} \code{lean_vm/src/cpu/mod.rs:568-570, 619-620},
  \code{flock/src/zerocheck.rs:389-390}, \code{lean_vm/src/hash_flock.rs:8-9} against
  \code{lean_vm/src/tables.rs:904-907}. \emph{Whose:} the Rust (comments). \emph{Asked:} a fix
  (\code{R112}).

\item \textbf{A stale comment on level-0 grinding.} It says level 0 grinds 0 bits in the
  production profile; the production derivation and its test set every level to 17. \emph{Where:}
  \code{pcs/src/whir.rs:1364-1365} (prover) and \code{:1771-1772} (verifier) against
  \code{pcs/src/whir_config.rs:60, 922, 1027} and \code{pcs/src/whir.rs:2485}. \emph{Whose:} the
  Rust (comments). \emph{Asked:} a fix (\code{R116}; the lines here were re-read at the pin and
  differ from the register's).

\item \textbf{A stale comment on the size of a WHIR Merkle leaf.} It says a leaf is
  \code{num_interleaved * 16} bytes; the code hashes $24$ bytes a lane. \emph{Where:}
  \code{pcs/src/whir.rs:394-395} against \code{:463}. \emph{Whose:} the Rust (a comment).
  \emph{Asked:} a fix (\code{R118}).
\end{enumerate}
\end{sloppypar}

\paragraph{What must not be reported.} The status file's finding that the Python verifier omits
the caps (its F9) is false at the pin: the Python's \code{build_layout} checks
$16 ≤ κ_{\mathrm{mem}} ≤ 32$, every height at most 32, $τ_{\mathrm{BLAKE2S}} ≥ 3$ and a
power-of-two program (\code{py:856-864}, called at \code{py:1378}; \code{log2_strict} at
\code{py:252-254}), and the Rust itself says the other two verifiers reject there
(\code{lean_vm/src/cpu/mod.rs:164-165}). It is to be deleted from the status, with the earlier
reviews that cite it, and nothing about it is to go upstream (\code{R38}). The review's findings
against the blueprint are leanerVM's own and are not leanVM matters either.
